% ============================================================
% Chapter: Information Completeness — Embedded Chapter
% Based on Anti-Hallucination Methodology v1.2 (老王审批 2026-06-17)
% Insert between §5 (Five-Anchor Audit) and §6 (False Positive Remediation)
% ============================================================

\section{Information Completeness: A Root Cause of Attribution Hallucination}
\label{sec:info-completeness}

The detection and audit mechanisms described in Sections~3--5 address threats 
originating from \textit{external} data manipulation. However, a second class of 
threat operates through a different mechanism: the model's own inference engine 
generating plausible but causally incorrect attributions when critical information 
is missing. This section formalizes the relationship between information gaps and 
attribution hallucination, presents a dual-perspective correction methodology, and 
integrates the resulting audit capability into the Five-Anchor framework.

\subsection{Problem Definition: The Info-Gap Hypothesis}

We define \textbf{attribution hallucination} as the phenomenon where a model 
generates a causal explanation that appears internally consistent but contains 
at least one link unsupported by verifiable evidence. Our central hypothesis is:

\begin{equation}
\label{eq:attribution}
H_{\text{attribution}} \Longleftrightarrow \text{info-gap} \land \text{strong-inference-closure}
\end{equation}

That is, attribution hallucination occurs exactly when two conditions are 
simultaneously satisfied: (1) an information gap exists---the model lacks access 
to facts necessary for sound causal attribution, and (2) the model possesses a 
strong inference engine that automatically fills the gap with ``plausible'' 
causal links derived from statistical patterns rather than empirical evidence.

\begin{figure}[ht]
\centering
\begin{tabular}{p{0.45\textwidth} p{0.45\textwidth}}
\toprule
\textbf{Without Gap} & \textbf{With Gap} \\
\midrule
$\text{Event}_A \xrightarrow{\text{Verified}} \text{Event}_B$ 
& $\text{Event}_A \xrightarrow{\text{Inferred}} \text{Motive}_X \xrightarrow{\text{Inferred}} \text{Event}_B$ \\
\cmidrule(lr){1-2}
All causal links independently verifiable 
& At least one link is a statistical bridge \\
\cmidrule(lr){1-2}
Low hallucination risk 
& High hallucination risk \\
\bottomrule
\end{tabular}
\caption{Info-gap as the structural precondition for attribution hallucination. 
The right column shows an inferred motive inserted to bridge a causal gap.}
\label{fig:info-gap}
\end{figure}

\subsubsection{Causal Link Taxonomy}

For audit purposes, we classify every link in a causal attribution chain into 
three types:

\begin{table}[ht]
\centering
\begin{tabular}{clp{0.5\textwidth}}
\toprule
\textbf{Type} & \textbf{Name} & \textbf{Definition} \\
\midrule
A & Directly Verifiable & Both endpoints independently confirmable via external sources \\
B & Partially Verifiable & One endpoint independently confirmable; the other inferred \\
C & Unverifiable Bridge & Neither endpoint independently confirmable; purely statistical \\
\bottomrule
\end{tabular}
\caption{Causal link classification for attribution audit.}
\label{tab:link-types}
\end{table}

\textbf{Trigger Rule:} Any Type~C link in a final attribution triggers 
\texttt{d2.5\_c\_gap\_trigger}, reducing attribution confidence proportionally 
to the ratio of Type~C links in the total causal chain.

\subsubsection{A Counter-Intuitive Prediction}

A distinguishing prediction of this framework is that \textbf{stronger models 
produce more attribution hallucination in info-gap scenarios} (see 
\S\ref{sec:experiments} for experimental design and predictions):

\begin{enumerate}
  \item Stronger model $\Rightarrow$ stronger inference engine $\Rightarrow$ faster inference-closure
  \item With the info-gap held constant, faster closure $\Rightarrow$ more confident but less detectable hallucination
  \item Weaker models are naturally safer in this dimension: insufficient inference capability prevents gap-filling
\end{enumerate}

This runs counter to the common assumption that model scaling universally 
improves reliability, and provides a falsifiable prediction for empirical testing.

\subsection{Methodology: Dual-Perspective Attribution Correction}

We present a five-stage protocol that combines two complementary perspectives: 
\textit{progressive gap exposure} (identifying where inference has substituted 
for evidence) and \textit{incentive traceback} (tracing stakeholder incentives 
to verify or refute inferred causal links).

\subsubsection{Stage 1: Progressive Gap Exposure}

Strong inference engines produce ``closed'' causal narratives where inferential 
jumps are indistinguishable from evidence-based links. The solution is not to 
suppress inference, but to systematically expose the gaps through multi-turn 
open-ended probing:

\begin{enumerate}
  \item For a given attribution claim $C$, identify all unverified premises 
        $\{p_1, p_2, \ldots, p_n\}$ that $C$ depends on
  \item For each premise $p_i$, ask: ``What evidence would falsify this premise?''
  \item If no falsification condition can be specified, mark the link as Type~C
  \item Iterate: each exposed gap may reveal sub-gaps at finer granularity
  \item Terminate when the gap set is empty or all remaining gaps are unqueryable
\end{enumerate}

This protocol is analogous to counterexample search in interactive proof 
systems, adapted for the non-automated context of human-AI collaborative 
attribution.

\subsubsection{Stage 2: Information Source Completion}

Once gaps are identified, the model performs targeted external retrieval for 
each gap. Unlike generic RAG, which retrieves broadly, this step retrieves 
specifically to fill identified causal gaps:

\begin{itemize}
  \item For each gap $g_i$, issue a directed query to external knowledge sources
  \item Verify that retrieved information directly addresses the identified gap
  \item Re-classify links after retrieval: Type~C $\rightarrow$ Type~B or Type~A
\end{itemize}

The principle: hallucination grows in information voids. Models fill 
empty spaces automatically, and the fill material is ``plausibility'' rather 
than ``truth.'' The solution is not to suppress inference, but to actively 
complete the information landscape---once the gaps are filled with evidence, 
inference no longer needs to bridge them.

\subsubsection{Stage 3: Incentive Chain Traceback}

For attribution claims involving multiple stakeholders, we trace the incentive 
structure of each stakeholder. The technical approach formalizes two 
co-isomorphic priors:

\begin{description}
  \item[Short-Chain Prior (Ockham)] Shorter causal explanations are preferred 
        over longer ones, all else being equal.
  \item[Internal-Incentive Prior] Internal incentive explanations are preferred 
        over external coercion explanations, all else being equal.
\end{description}

These two priors are structurally isomorphic: both prefer explanations with 
fewer unobserved intermediaries. The combined rule---the \textbf{Self-Interest 
Prior}---states: \textit{exhaust internal incentive explanations before 
invoking external intervention.}

\subsubsection{Stage 4: Self-Interest Prior Convergence}

Given the retrieved information sources $S$ and the stakeholder incentive map 
$I$, the correction algorithm converges on the most parsimonious explanation:

\begin{enumerate}
  \item Generate the shortest causal chain consistent with all verified sources $S$
  \item Identify the most sufficient internal incentive explanation from $I$
  \item If the internal incentive explanation is sufficient, adopt it
  \item Otherwise, adopt the shortest-chain explanation
  \item If neither is sufficient, mark the attribution as \texttt{UNCERTAIN}
\end{enumerate}

\subsubsection{Stage 5: Attribution Verification}

After correction, verify that the revised attribution chain contains no 
unacknowledged inference bridges. If any Type~C links remain:

\begin{itemize}
  \item Mark them explicitly with \texttt{[INFERRED]} tags
  \item Compute confidence as $1 - (n_C / n_{\text{total}})$, where $n_C$ is the number of Type~C links
  \item Flag for human review if confidence falls below threshold
\end{itemize}

\subsubsection{Algorithm}

\begin{small}
\begin{verbatim}
function dual_perspective_correction(claim, domain):
    // Stage 1: Progressive Gap Exposure
    gaps = open_ended_probe(claim)
    if gaps.empty() or gaps.all_unqueryable():
        return {verdict: GAP_BLOCKED}

    // Stage 2: Information Source Completion
    sources = [external_retrieval(g) for g in gaps]

    // Stage 3: Incentive Chain Traceback
    stakeholders = identify_stakeholders(domain)
    incentive_map = {s: compute_incentive(s, claim) for s in stakeholders}

    // Stage 4: Self-Interest Prior Convergence
    // co-isomorphic: both priors select the most parsimonious explanation
    short_chain = shortest_causal_explanation(sources)
    internal = most_sufficient_internal_incentive(incentive_map)
    
    correction = internal if internal.sufficient else short_chain

    // Stage 5: Verification
    if correction == null or correction.depends_on(\"inference_bridge\"):
        inference_bridges = list_inference_bridges(claim)
        return {verdict: UNCERTAIN,
                gaps: gaps,
                gap_types: classify_gaps(gaps),
                inference_bridges: inference_bridges,
                confidence: len(inference_bridges) / total_links(claim)}

    return {verdict: CORRECTED, correction: correction}
\end{verbatim}
\end{small}

\subsection{Degradation Spectrum}

The methodology's effectiveness degrades along a spectrum determined by the 
nature of the information gap:

\begin{center}
\begin{tabular}{lp{0.55\textwidth}}
\toprule
\textbf{Stage} & \textbf{Description} \\
\midrule
Fully Queryable & Complete correction possible; all gaps can be exposed and filled \\
Partially Queryable & Partial correction: gaps can be identified but not all filled; remaining uncertainty is quantified \\
Unqueryable & Degraded: can only flag global uncertainty without offering alternative attribution \\
Hidden Gap & Worst case: the info-gap itself is masked by inference closure; methodology fails silently \\
\bottomrule
\end{tabular}
\end{center}

The \texttt{GAP\_BLOCKED} return in the algorithm covers the latter two stages. 
The first two stages are differentiated through the \texttt{UNCERTAIN} return's 
quantified confidence value and gap-type classification.

\subsection{Comparison with Existing Approaches}

\begin{table}[ht]
\centering
\begin{tabular}{lp{0.35\textwidth}p{0.4\textwidth}}
\toprule
\textbf{Approach} & \textbf{Mechanism} & \textbf{Limitation} \\
\midrule
RLHF/RLAIF & Preference labeling suppresses undesired outputs & Behavioral constraint; treats symptoms, not causes \\
RAG & Retrieval fills knowledge gaps & Fills factual gaps only; does not address causal attribution chains \\
Multi-Agent Voting & Consensus via debate or majority & Shared info-gap $\rightarrow$ shared hallucination is a blind spot \\
Debate + Human Judge & AI debate with human arbitration & Approaches our method but more adversarial; we prioritize collaborative correction \\
Herrera (2026) Architecture & Engineering pipeline for hallucination governance & Their identified ``next step'' is our existing architecture \cite{herrera2026} \\
\midrule
\textbf{This Work} & Gap exposure + incentive traceback + Self-Interest Prior & \textbf{Cognitive audit}: addresses root cause rather than output behavior \\
\bottomrule
\end{tabular}
\caption{Comparison with existing hallucination mitigation approaches.}
\label{tab:comparison}
\end{table}

While each individual component (gap detection, retrieval, incentive modeling) 
has academic precedent, the integration of all components into a unified 
cognitive audit architecture is novel. The key differentiator is treating 
hallucination not as an output quality problem but as a structural consequence of 
the inference engine operating with incomplete information.

\subsection{Framework Integration}

The methodology integrates into the existing Five-Anchor framework at six 
specific touchpoints, strengthening the audit capability without adding new 
anchors or disrupting existing structures:

\subsubsection{D2: Logic Anchor --- Attribution Audit}

\begin{itemize}
  \item \textbf{d2.4\_info\_gap\_detection}: Add a column to the logic audit 
        matrix classifying each causal link as Type A/B/C
  \item \textbf{d2.5\_c\_gap\_trigger}: Any Type~C link triggers automatic 
        confidence downgrade for the attribution verdict
\end{itemize}

\subsubsection{L5: Causal Attribution Layer --- New Sub-layers}

\begin{itemize}
  \item \textbf{L5b\_interest\_chain}: Incentive traceback sub-module that maps 
        stakeholder incentives for any multi-party attribution claim
  \item \textbf{L5c\_self\_interest\_prior}: Self-Interest Prior convergence 
        module that selects the most parsimonious internal-incentive explanation
\end{itemize}

The updated detection chain structure becomes:

\begin{center}
\begin{tabular}{l}
L1: Basic Fact Matching \\
L2: Semantic Consistency \\
L3: Logical Closure \\
L4: Statistical Bias Detection \\
L5: Causal Attribution Audit \\
\quad L5a: Surface Causality Check \\
\quad L5b: Incentive Chain Traceback \hfill \textit{[new]} \\
\quad L5c: Self-Interest Prior Convergence \hfill \textit{[new]} \\
L6: Counter-Cognitive-Colonization Audit \\
\end{tabular}
\end{center}

\subsubsection{D5: Methodology Anchor --- New Rules}

\begin{itemize}
  \item \textbf{d5.5\_prior\_rule}: Self-Interest Prior as a cross-anchor 
        convergence rule---when multiple anchors produce conflicting verdicts 
        on attribution claims, the prior guides resolution
  \item \textbf{d5.6\_info\_threshold}: Maximum allowable Type~C link ratio in 
        any core causal chain, initially set conservatively at $\leq 1$. This 
        parameter is flagged for dynamic adjustment via FML calibration data
\end{itemize}

\textbf{Implicit Assumption Correction.} The info-gap hypothesis directly 
challenges an unstated assumption in the methodology anchor: that sufficiently 
good methodology can compensate for missing information. Our analysis shows the 
opposite---stronger methodology accelerates gap-closure, making hallucination 
more confident and less detectable. Therefore, \texttt{d5.6\_info\_threshold} 
functions as a brake mechanism rather than an optimization target.

\subsubsection{FML: Failure Memory Layer}

Info-gap patterns are stored in the Failure Memory Layer for future 
pre-triggering. When a new attribution task matches a stored gap pattern, the 
system automatically reduces inference speed and increases verification 
requirements before generating attribution.

\subsubsection{Tier IV: Cross-Cluster Detection}

A new coroutine \texttt{tier4\_info\_gap\_coroutine} detects whether apparent 
multi-source consensus actually originates from a shared information gap---a 
pattern where multiple sources independently confirm the same causal narrative 
because they share the same missing information, not because the narrative is 
true.

\subsubsection{D6: Language-View Layer}

Attribution audit tags are added to the language-view annotation schema:
\texttt{[GAP]} for identified information gaps, \texttt{[INFERRED]} for 
inference-bridge links, and \texttt{[VERIFIED]} for independently confirmed 
causal links.

\subsection{Experimental Design and Predictions}
\label{sec:experiments}

The experimental methodology follows directly from the info-gap hypothesis 
(\S\ref{eq:attribution}): construct a fixed set of info-gap scenarios 
(missing key stakeholder motivation information), expose multiple models of 
varying capability (7B/13B/70B/DeepSeek-R1) to identical scenarios, and 
measure the ratio of Type~C links in causal attributions alongside independent 
fact-checking accuracy. A control condition uses the same scenarios with 
info-gaps filled. This design isolates the effect of inference capability on 
attribution quality independent of scenario difficulty.

The framework generates specific, falsifiable predictions suitable for 
empirical validation:

\begin{enumerate}
  \item \textbf{Model Scaling $\times$ Hallucination Rate.} In controlled 
        info-gap scenarios, larger models will exhibit a higher Type~C link 
        ratio than smaller models, while simultaneously expressing higher 
        confidence in their attributions. This creates a ``confidence-hazard 
        inversion'' where stronger models are more dangerous for attribution 
        tasks with incomplete information.

  \item \textbf{Gap-Filling Efficacy.} Applying the five-stage protocol will 
        reduce Type~C link ratios from a baseline (estimated 20--40\% for strong 
        models in info-gap conditions) to below 5\%, with the majority of 
        remaining Type~C links classified as Known Unknowns.

  \item \textbf{Self-Interest Prior Accuracy.} In multi-stakeholder scenarios 
        where ground-truth incentives are independently verifiable, the 
        Self-Interest Prior will correctly identify the dominant causal factor 
        in $\geq 80\%$ of cases, outperforming both pure chain-of-thought 
        reasoning and external-intervention-first heuristics.

  \item \textbf{Shared Info-Gap Detection.} The Tier~IV gap-coroutine will 
        identify shared info-gap consensus patterns with a false positive rate 
        below 15\% in controlled multi-source scenarios.
\end{enumerate}

\subsection{Applicability Boundaries}

The methodology explicitly does not apply to all hallucination types. Its 
effectiveness is limited to \textbf{attribution hallucination}---cases where the 
model constructs causal narratives with information gaps. It is not designed for:

\begin{itemize}
  \item Factual hallucination (missing knowledge): better addressed by retrieval 
        augmentation and confidence calibration
  \item Statistical bias hallucination: requires rare-signal amplification methods
  \item Semantic ambiguity hallucination: requires disambiguation anchoring
  \item Adversarial injection: requires security-hardening mechanisms
  \item Value-alignment hallucination: requires broader calibration framework 
        beyond any single methodology
\end{itemize}

Within its domain, however, the methodology addresses a failure mode that 
existing approaches systematically miss: the \textit{invisible} hallucination 
that appears structurally sound but contains unverifiable causal links.
